\documentclass[a4paper,11pt]{ctexart}
\usepackage{xcolor}
\usepackage[top=2cm, bottom=2cm, left=2cm, right=2cm]{geometry}
\usepackage{array}
\usepackage{hyperref}
\usepackage{booktabs}
\usepackage{amsmath}
\usepackage{amssymb}
\hypersetup{colorlinks=true,linkcolor=blue!70!black}
\linespread{1.35}

\newcommand{\done}{\textcolor{green!50!black}{\textbf{【已有初稿】}}}
\newcommand{\draft}{\textcolor{orange!80!black}{\textbf{【待写】}}}

\title{\textcolor{blue!70!black}{\Huge\textbf{金融学模块 · 完整课程大纲}}}
\subtitle{\textcolor{gray}{三卷本 · AP/ A-Level 至大一金融专业先修水平}}
\author{}
\date{}

\begin{document}
\maketitle
\thispagestyle{empty}

\noindent\textbf{关于课程设计：}

\begin{quote}
本模块不是理财科普，而是一套严肃的金融学先修课程。定位对标：AP Macroeconomics (Financial Sector) + AP Microeconomics (Market Structures) + 部分CFA Level I 核心内容。三卷共24课，每课按统一结构组织：\\
\textbf{情景引入} → \textbf{定性分析} → \textbf{定量分析} → \textbf{总结} → \textbf{讨论题} → \textbf{课后练习} → \textbf{学海拾遗（学术花絮）}
\end{quote}

\noindent\textbf{与同项目其他模块的关系：}
\begin{itemize}
  \item \textit{《经济小故事》}（社科系列，G4-G7）——完全重叠，本模块不引用、不依赖
  \item E04-intro/（入门故事，G5-G8）——提供直觉铺垫，但不作为先修要求
  \item \textbf{本三卷本面向已完成 Algebra II 和基础概率的学生}
\end{itemize}

\newpage

% ============================================================
\section*{\textcolor{blue!70!black}{第一卷\quad 金融学基础（Foundations of Finance）}}
% ============================================================

\textbf{内容：} 货币的时间价值、债券与固定收益、股票估值、现代投资组合理论、CAPM、市场效率、央行与货币政策。\vspace{0.3cm}

\begin{enumerate}

\item[\textbf{第1课}] \textbf{货币的时间价值——一切定价的起点} \done

\textbf{情景：} 一个北京家庭需要为10年后的大学教育金做规划。现有存款50万，目标是100万。怎么实现？

\textbf{定性：} 为什么今天的一块钱比明天的一块钱值钱？机会成本、通胀、不确定性三个根源。

\textbf{定量：} $FV = PV(1+r)^n$、$PV = FV/(1+r)^n$、72法则、连续复利 $FV = PV \cdot e^{rt}$、EAR vs APR 的转换。

\textbf{讨论：} "72法则在低利率时误差大"——在高利率（>20\%）和低利率（<2\%）时用69.3法则更准。为什么？

\textbf{练习：} 给定一组现金流序列，计算NPV和IRR。分析两个互斥项目的选择。

\textbf{学海拾遗：} 贴现率的哲学——为什么经济学用折现而不用"上浮"？Ramsey方程与社会贴现率的伦理争论。

\item[\textbf{第2课}] \textbf{债券定价与利率风险} \done

\textbf{情景：} 2020年买入10年期国债的投资者，到2024年发现债券价格跌了——为什么央行加息会让已发行的债券贬值？

\textbf{定性：} 票面利率 vs 到期收益率 vs 即期利率。折价/平价/溢价债券。利率与价格的跷跷板关系。

\textbf{定量：} 
\[
P = \sum_{t=1}^{n} \frac{C}{(1+y)^t} + \frac{F}{(1+y)^n}
\]
久期（Macaulay Duration）与修正久期：$\Delta P/P \approx -D^* \cdot \Delta y$。凸性校正。

\textbf{讨论：} 为什么长期债券比短期债券对利率变动更敏感？零息债券的久期为什么等于期限？

\textbf{练习：} 计算三只不同票面利率、不同期限债券的久期和凸性。给定利率上升200bp，预估价格变化并与精确计算对比。

\textbf{学海拾遗：} 免疫策略（Immunization）——养老金基金如何用久期匹配来锁定未来负债。为什么2023年硅谷银行（SVB）在持有"最安全"国债的情况下还是破产了？久期错配的致命后果。

\item[\textbf{第3课}] \textbf{收益率曲线与期限结构} \draft

\textbf{情景：} 2024年8月，中国10年期国债收益率2.1\%，2年期1.5\%——长端高于短端（正常）。但2022年出现过2年期高于10年期的"倒挂"——被视为经济衰退的先兆。

\textbf{定性：} 为什么长期利率一般高于短期？流动性偏好理论、预期理论、市场分割理论。

\textbf{定量：} 即期利率曲线拟合（Bootstrapping）、远期利率计算：
\[
(1+f_{t,T})^{T-t} = \frac{(1+y_T)^T}{(1+y_t)^t}
\]

\textbf{讨论：} 收益率曲线倒挂为什么被当作经济衰退信号？过去50年每次美国衰退前都出现了倒挂——但它不一定"导致"衰退，而是市场对未来的预期集中表达。

\textbf{练习：} 给出一组国债价格，用 bootstrapping 推导即期利率曲线和远期利率曲线。根据远期利率推断市场对未来利率的预期。

\textbf{学海拾遗：} 中国收益率曲线的特殊性——国债市场分割（商业银行持有至到期为主、交易盘占比低）导致曲线形态与美国有本质不同。信用债 vs 利率债的利差（Credit Spread）。

\item[\textbf{第4课}] \textbf{股票估值——从Gordon模型到DCF} \done

\textbf{情景：} 2020年末茅台市盈率超过60倍——年增速15\%的预期能支撑这个估值吗？2024年茅台市盈率降到30倍——发生了什么？

\textbf{定性：} 股票价值 = 未来现金流折现之和。增长的"乘法器效应"——小幅增长预期变化导致估值大幅波动。

\textbf{定量：} Gordon增长模型 $P_0 = D_1/(r-g)$、多阶段增长模型（H模型）、P/E比率与 Gordon 模型的关联 $P/E = (1-b)/(r-ROE \times b)$、DCF全流程（自由现金流预测→WACC计算→终值估计）。

\textbf{讨论：} PEG比率（P/E÷增长率）的分母用过去增长还是未来增长？DCF模型中最敏感的参数是什么——一个参数的微小变化如何导致估值翻倍或腰斩？

\textbf{练习：} 给出一家公司的财务数据，建立三阶段DCF模型。做敏感性分析（当WACC变化±1\%，终值增长率变化±0.5\%时，估值怎么变）。

\textbf{学海拾遗：} DCF的哲学困境——它假设未来可预测。亚马逊连续亏损十几年，如果1999年用DCF估值，大概率给出"卖出"评级。但持有到今天涨了1000倍。DCF在高增长、不盈利、无形资产密集的公司面前几乎完全失效——替代方法：实物期权定价。

\item[\textbf{第5课}] \textbf{现代投资组合理论（MPT）} \done

\textbf{情景：} "不要把所有鸡蛋放在一个篮子里"——但多少只鸡蛋才算够？15只？30只？为什么不是越多越好？

\textbf{定性：} 分散化的数学本质：协方差和相关系数。系统性风险 vs 非系统性风险。

\textbf{定量：} 
\[
\sigma_p^2 = \sum_{i=1}^{n} w_i^2 \sigma_i^2 + \sum_{i=1}^{n} \sum_{j \neq i} w_i w_j \sigma_i \sigma_j \rho_{ij}
\]
有效前沿（Efficient Frontier）的推导、最小方差组合、Market Portfolio。

\textbf{讨论：} 当 $n$ 从1增加到50时，组合风险如何下降？实证数据显示15-20只股票可以消除约90\%的非系统性风险——余下的10\%是系统性风险，无法通过分散化消除。

\textbf{练习：} 给定5只股票的历史收益率和协方差矩阵，用Excel或Python求解有效前沿。找出最小方差组合和最大夏普比率组合。

\textbf{学海拾遗：} Markowitz的均值-方差框架发表于1952年——他在博士论文答辩时，Milton Friedman 说"这不是经济学，这是数学"。62年后Markowitz获得了诺贝尔奖。有效前沿的真实表现：由于参数估计误差（输入敏感性问题），样本外有效前沿经常退化——Michaud（1989）称之为"估计误差最大化器"。

\item[\textbf{第6课}] \textbf{资本资产定价模型（CAPM）与套利定价理论（APT）} \done

\textbf{情景：} 为什么高$\beta$股票在牛市中暴涨，在熊市中暴跌？市场"奖励"承担系统性风险的投资者——但如何度量这个"奖励"的幅度？

\textbf{定性：} CAPM的假设体系（有效市场、无限借贷、共同预期）——这些假设在现实中全都不成立，但模型为什么还有用？

\textbf{定量：} 
\[
E(R_i) = R_f + \beta_i [E(R_m) - R_f]
\]
$\beta$的估计方法（历史回归 vs 调整$\beta$）、证券市场线（SML）、詹森$\alpha$。

\textbf{讨论：} CAPM受到大量实证挑战——Roll的批评（市场组合不可观测）、Banz的小公司效应、Stattman的账面市值比效应、Fama-French三因子模型凭什么比CAPM更好？

\textbf{练习：} 对一组A股股票做$\beta$计算（至少3年日收益率回归）。在中国市场，CAPM是否能解释截面收益率差异？计算每只股票的$\alpha$，判断哪些"跑赢"了市场。

\textbf{学海拾遗：} Fama-French五因子模型（市场、规模、价值、盈利、投资）和Carhart四因子模型（+动量）。在中国A股市场，小盘因子和动量因子的表现与美股截然不同——市场微观结构差异决定了因子在不同市场的有效性。

\item[\textbf{第7课}] \textbf{有效市场假说与行为金融} \draft

\textbf{情景：} 1987年10月19日"黑色星期一"，道琼斯指数一天暴跌22.6\%——没有明显利空。如果市场是有效的，这么大的跌幅必须在当天有对应量的新闻——但那天没有。

\textbf{定性：} EMH三种形式的逻辑和检验。随机游走假说。事件研究法。

\textbf{定量：} 通过事件研究法检验半强式有效——盈余公告后的股价漂移（PEAD，Ball \& Brown 1968）。

\textbf{讨论：} 为什么主动基金经理长期跑不赢指数？Fama vs Shiller——同一个诺贝尔奖分给两个观点完全相反的人。市场有效假说不是一个"是或否"的问题——它是市场效率的"基准"，效率程度因市场、时间、投资者结构而异。

\textbf{练习：} 选一个A股市场的重大事件（如降准、贸易战升级），用事件研究法分析股价在事件前后的反应速度和准确性。

\textbf{学海拾遗：} 行为金融学的关键发现——前景理论（Kahneman-Tversky）、过度反应与反应不足（De Bondt-Thaler）、处置效应（Shefrin-Statenman）、本地偏好（Home Bias）。Richard Thaler的"心理账户"——为什么人们对"赚来的钱"和"意外之财"采取不同的风险态度。

\item[\textbf{第8课}] \textbf{中央银行与货币政策} \done

\textbf{情景：} 2022-2023年，美联储一年加息11次（0.25\%→5.5\%），而中国央行同步降息。两国央行为什么做出了截然不同的选择？谁的决策是对的？

\textbf{定性：} 央行的三个使命（价格稳定/充分就业/金融稳定）、货币政策的传导机制（利率渠道/信贷渠道/资产价格渠道/汇率渠道）。

\textbf{定量：} 货币乘数 $m = 1/rr$、Taylor Rule $i = r^* + \pi + 0.5(\pi - \pi^*) + 0.5(y - y^*)$、数量方程式 $MV = PY$。

\textbf{讨论：} 中国的货币政策传导效率为什么低于美国？——利率市场化尚未完全完成、国企/地方政府对利率不敏感、商业银行的风险偏好约束。

\textbf{练习：} 用Taylor Rule估算当前中国的政策利率是否"适当"。分析2024年9月中国央行降准50bp对货币乘数和M2的量化影响。

\textbf{学海拾遗：} 现代货币理论（MMT）——政府可以无限印钞而不违约吗？日本国债/GDP超过260\%但从未违约，为什么？"货币融资"（直接印钞给财政）vs "债务融资"——日本央行YCC（收益率曲线控制）的实践与退出困境。

\end{enumerate}

\newpage

% ============================================================
\section*{\textcolor{blue!70!black}{第二卷\quad 金融衍生品与风险管理（Derivatives \& Risk Management）}}
% ============================================================

\textbf{说明：} 衍生品是金融学的硬核部分——定价逻辑极其精致、但也因其复杂性导致了多次重大灾难。本卷假设学生已掌握第一卷的定价基础。

\begin{enumerate}

\item[\textbf{第9课}] \textbf{远期合约与期货——锁定未来的价格} \draft

\textbf{情景：} 一个北京的面包商担心小麦价格上涨；与此同时一个河北农民担心小麦收获时价格下跌。两个人如何通过一份合约各取所需？

\textbf{定性：} 远期 vs 期货的区别（标准化/保证金/逐日盯市/交割方式）。套期保值（Hedging）与投机（Speculation）的本质区别。

\textbf{定量：} 远期价格定价：
\[
F = S_0 \times e^{(r+q)T}
\]
（$q$为便利收益率/存储成本）。期货基差（Basis）与基差风险。最优对冲比率 $h^* = \rho \times (\sigma_S / \sigma_F)$。

\textbf{讨论：} 中国商品期货市场——大连商品交易所（大豆/玉米）、上海期货交易所（铜/螺纹钢）、郑州商品交易所（白糖/PTA）——为什么中国是全球最大的商品期货市场之一？

\textbf{练习：} 一家中国航空公司担心未来油价上涨，如何使用原油期货对冲？计算合约数量。比较期货对冲与直接买入原油现货的优劣。

\textbf{学海拾遗：} 中航油事件（2004）——中国航油（新加坡）在石油期货上投机巨亏5.5亿美元。套保还是投机？区分二者的关键在于：交易方向是否与你的现货风险敞口相反。

\item[\textbf{第10课}] \textbf{互换（Swaps）——利率互换与信用违约互换} \draft

\textbf{情景：} 一家中国企业借了浮动利率美元贷（LIBOR+200bp），但预期美联储会加息。一家中国银行持有大量固定利率国债，但预期利率会下行。它们如何互换各自的风险？

\textbf{定性：} 利率互换（IRS）的机制——固定端vs浮动端。比较优势理论——为什么通过互换双方可以获得更好的融资条件？

\textbf{定量：} IRS定价——固定端利率等于使互换净现值为零的固定比率：
\[
\text{固定利率} = \frac{1 - DF_n}{\sum_{t=1}^{n} DF_t}
\]
信用违约互换（CDS）的基础结构——买方定期支付保费、卖方在违约发生时赔付。

\textbf{讨论：} CDS是"保险"还是"赌博"？——如果你不持有标的债券却买入CDS（裸CDS）——你是在赌它违约。2008年AIG就是因为卖出了大量的CDS却未计提足够保证金而倒闭的。1.5万亿美元的CDS市场——其中80\%以上是"裸"的。

\textbf{练习：} 给定一组LIBOR远期利率，计算一个3年期IRS的固定利率。分析企业选择固定利率vs浮动利率贷款的决策因素。

\textbf{学海拾遗：} CDS大狂潮——2000-2007年CDS名义本金从1万亿美元飙升至62万亿美元（超过了全球GDP）。没有任何集中清算平台、没有任何保证金要求、没有任何监管——它是"影子银行"体系的核心组带。2008年之后CDS市场被严格监管，但到2024年名义本金又回到了约10万亿美元。

\item[\textbf{第11课}] \textbf{期权基础——权利与义务的分离} \draft

\textbf{情景：} 2021年Archegos Capital Management通过总收益互换（TRS）杠杆买入了大量中概股和媒体股。当股价暴跌时，Archegos无法追加保证金，造成全球银行损失超过100亿美元。期权的杠杆从何而来？

\textbf{定性：} 看涨/看跌期权的基本结构。欧式vs美式。实值/平值/虚值。

\textbf{定量：} 到期时的盈亏图：
\begin{align*}
\text{Call: } & \max(S_T - K, 0) \\
\text{Put: } & \max(K - S_T, 0)
\end{align*}
看涨-看跌平价关系（Put-Call Parity）：
\[
C + PV(K) = P + S_0
\]

\textbf{讨论：} 期权为什么是"非线性"的？——股票的$\Delta$固定为1，而期权的$\Delta$是动态变化的。这让期权能构造出任何你想得到的收益形态。期权复制的思想——BS公式的数学根源。

\textbf{练习：} 画（或编写绘图代码画出）跨式组合（Straddle）、勒式组合（Strangle）、蝶式价差（Butterfly）的盈亏图。分析一个"保护性看跌"（Protective Put）策略的最大损失和盈亏平衡点。

\textbf{学海拾遗：} 期权市场的一个关键数字——未平仓量/成交量比率。90\%的期权在到期日变为虚值。购买期权是一种被称为"高概率亏损、低概率暴赚"的博弈。专业交易员更多是卖出期权（赚取时间价值）而不是买入期权。

\item[\textbf{第12课}] \textbf{期权定价——Black-Scholes模型} \draft

\textbf{情景：} Myron Scholes和Fischer Black在1969-1973年间推导出了期权定价公式——它如此优美，以至于Scholes和Merton（不是Black，他1995年去世了）在1997年获得了诺贝尔奖。但获奖两年后，他们参与的对冲基金LTCM破产了……

\textbf{定性：} Black-Scholes的五个输入——股价、执行价、到期时间、波动率、无风险利率。唯一不可观测的输入是波动率——但它决定了所有价格。

\textbf{定量：}
\[
C = S_0 N(d_1) - K e^{-rT} N(d_2)
\]
\[
d_1 = \frac{\ln(S_0/K) + (r + \sigma^2/2)T}{\sigma\sqrt{T}}, \quad d_2 = d_1 - \sigma\sqrt{T}
\]
风险中性定价的直觉——为什么BS公式中不含股票的预期收益？

\textbf{讨论：} 隐含波动率微笑（Volatility Smile）——1987年股灾后，市场意识到BS假设的对数正态分布低估了极端事件概率。波动率曲面（Vol Surface）是市场对BS公式的"补丁"。

\textbf{练习：} 给定BS参数，用Python（或计算器）计算期权价格。计算期权的Greeks（Delta, Gamma, Theta, Vega, Rho）。用隐含波动率反推市场预期。

\textbf{学海拾遗：} LTCM（1998年）——两个诺贝尔奖得主管理的对冲基金怎么破产的？他们的模型假设"历史的波动率就是未来的波动率"——1998年俄罗斯债务违约引发的连锁反应是模型中没有的"10个标准差事件"。LTCM的倒闭不是因为模型错了，而是因为他们用50倍杠杆放大了模型的风险暴露——当市场不再"正常"时，清算所要求追加的保证金超过了他们的剩余资本。

\item[\textbf{第13课}] \textbf{风险管理——VaR、压力测试与久期匹配} \draft

\textbf{情景：} 1994年宝洁公司（P&G）的财务主管签署了一份复杂的利率互换协议——他以为利率不会大幅上升。结果美联储加息250bp，P&G亏了1.57亿美元。为什么会签这样的协议？因为风险被隐藏在了复杂的衍生品结构中。

\textbf{定性：} 风险管理的四步：识别→度量→管理→监测。市场风险/信用风险/流动性风险/操作风险/模型风险的区别。

\textbf{定量：} VaR（Value at Risk）的计算方法——历史模拟法、方差-协方差法、蒙特卡洛模拟法。VaR的缺陷：它不告诉你"超过VaR的损失有多大"。
\[
\text{ES（Expected Shortfall）} = E[L | L > \text{VaR}]
\]
巴塞尔协议对商业银行资本金的要求。

\textbf{讨论：} VaR的致命缺陷——2008年金融危机前，华尔街各大银行的VaR模型显示"99\%安全的"。但这个模型根本不包括流动性风险和相关性突变——当所有资产同时暴跌（相关性趋近1），VaR完全失去了意义。

\textbf{练习：} 给定一个投资组合的历史收益率，计算95\%和99\%置信水平下的VaR和Expected Shortfall。构建一个简单的压力测试情景（如市场暴跌20\%、利率上升200bp）。

\textbf{学海拾遗：} 久期缺口管理——银行的资产负债管理（ALM）核心。如果一个银行的资产久期大于负债久期（常态），利率上升会使银行净值缩水。硅谷银行（2023）的资产久期约6年、负债久期约1年——当美联储加息500bp后，6.2\%的资产价值跌幅远远超过了银行的资本充足率。久期匹配不是"稳健"的选择——它是"稳健"的数学表达。

\item[\textbf{第14课}] \textbf{结构型产品与资产证券化} \draft

\textbf{情景：} 2007年，一只名为"合成CDO"的金融产品——它甚至不持有任何真实的房贷，只是对一篮子CDS的"再打包"——获得了AAA评级。为什么评级机构会给这种"再打包再打包再打包"的产品最高评级？

\textbf{定性：} 资产证券化的流程——贷款池 → SPV → 分级（Senior/Mezzanine/Equity）→ 销售。信用增级机制（超额担保/优先级结构/利差账户）。

\textbf{定量：} 违约相关性模型——Gaussian Copula（正态连接函数）——David Li在2000年发表的论文，被华尔街广泛用于CDO定价。$ \rho $ 的选择决定了整个定价结果。

\textbf{讨论：} Gaussian Copula被批评为"导致金融危机的数学公式"——它假设违约事件之间的相关性是恒定的（不随市场状态变化），而在危机中所有相关性都趋近于1。2008年后结构化金融被妖魔化——但证券化本身不是问题，问题在于劣质基础资产+过度分级+评级腐败+零透明度。

\textbf{练习：} 简化模拟一个CDO的分级结构——给定资产池违约概率和违约相关性，计算Senior/Mezzanine/Equity三个层级的预期损失和评级。

\textbf{学海拾遗：} 中国资产证券化市场——2024年存量规模约4万亿（vs 美国12万亿），以信贷ABS和RMBS为主。中国ABS市场的特殊问题：基础资产池的详情披露不充分、发起人自持比例不足、法律上的"真实出售"认定不清。另一个结构型产品——结构性存款（保本+衍生品）——在中国卖了超过10万亿，但很多投资者的理解停留在"保本高息"层面，完全不知道衍生品部分的"敲出"条款会让他们损失全部收益。

\end{enumerate}

\newpage

% ============================================================
\section*{\textcolor{blue!70!black}{第三卷\quad 金融创新与系统性风险（Innovation \& Systemic Risk）}}
% ============================================================

\textbf{说明：} 本卷将前两卷的工具用于理解金融史上最大的几次崩溃和最新的金融前沿。

\begin{enumerate}

\item[\textbf{第15课}] \textbf{影子银行——监管之外的大玩家} \draft

\textbf{情景：} 2019年中国包商银行被接管——原因是银行间市场"同业存单"的刚性兑付被打破。包商银行的同业存单规模是其存款的2倍——它是一家"影子银行化"的商业银行。

\textbf{定性：} 影子银行的定义（金融稳定理事会FSB：信贷中介活动中那些在常规银行体系之外的、存在监管套利的部分）。核心特征：期限转换（借短投长）、流动性转换、高杠杆、监管缺席。

\textbf{定量：} 中国影子银行的主要形式——银行理财（巅峰期约30万亿）、信托贷款、委托贷款、未贴现银行承兑汇票。2017年金融去杠杆启动后，影子银行规模从100万亿收缩到约60万亿。

\textbf{讨论：} 影子银行是"金融创新"还是"监管套利"？——通常是两者兼有。它们提高了信贷供给效率，但也创造了系统性风险。中国的"资管新规"（2018）如何尝试将影子银行纳入监管？

\textbf{练习：} 计算一个银行理财产品的"资金池"模式下——当底层资产违约率从2\%升至10\%时，银行需要补充多少资本金。

\textbf{学海拾遗：} 中国的"同业存单—同业理财—委外投资"循环——形成了一条独立于存贷体系的影子信贷供应链。2016-2017年的"债灾"正是因为这条供应链断裂造成的。美国的影子银行——货币市场基金、私募信贷基金、Repo市场——在2023年3月的硅谷银行危机中再次暴露了它的脆弱性。

\item[\textbf{第16课}] \textbf{2008年全球金融危机——深度解剖} \draft

\textbf{情景：} 2008年9月15日，雷曼兄弟——158年历史的投资银行——申请破产保护。全球金融体系在48小时内冻结。为什么一个美国投行的倒闭能导致全球性的大萧条以来最严重的危机？

\textbf{定性：} 危机的七步链：低利率环境→房贷泡沫→次贷证券化→评级腐败→杠杆堆积→房价转跌→多米诺坍塌。

\textbf{定量：} MBS构造、CDO分层、CDS保险的互联性。AIG的CDS敞口超过5000亿美元——它卖了保险却没有持有足够的资本。AIG的倒闭被美国政府以1800亿美元救助阻止——因为"大而不能倒"（Too Big to Fail）。

\textbf{讨论：} "大而不能倒"的道德风险——救助保住了金融体系，但让大型金融机构有了"亏了政府兜底"的预期。Dodd-Frank法案（2010）试图解决这个问题——但2023年硅谷银行倒闭说明问题并没有解决。更有力的解决方案：更高的资本金要求、生前遗嘱（Living Wills）、系统性风险附加费。

\textbf{练习：} 构建一个简化的资产负债表传染模型——如果A银行违约，哪些与A银行有同业往来的银行会受影响？（用矩阵写出银行间风险暴露的传播路径。）

\textbf{学海拾遗：} "格林斯潘对策"（Greenspan Put）——2000年互联网泡沫破裂后，美联储将利率从6.5\%降至1\%，催生了后来的房地产泡沫。2004-2006年美联储又快速加息至5.25\%（加息17次）——直接刺破了房地产泡沫。央行的每一次干预，都在为下一次危机埋下种子。"事后诸葛亮"很容易，但危机中和危机前的决策完全不同。

\item[\textbf{第17课}] \textbf{LTCM、长期资本管理与模型风险} \draft

\textbf{情景：} 1998年，LTCM——由两位诺贝尔奖得主（Scholes, Merton）和华尔街最优秀的交易员管理——在几周内从50亿美元资本亏到几乎为零。美联储紧急组织了14家银行注资36.5亿美元。诺贝尔奖得主的模型，为什么不管用了？

\textbf{定性：} LTCM的策略——收敛交易（Convergence Trading）：买入被低估的债券、卖出被高估的债券，赌价差回归历史均值。规模超过1000亿美元（50倍杠杆）。

\textbf{定量：} LTCM的模型假设——基于过去5年的历史波动率和相关性。但1998年8月俄罗斯违约后，全球避险情绪让所有相关性同时趋近于1——"分散化投资"的假设瞬间失效。市场流动性消失——LTCM无法以不亏损的价格平仓。

\textbf{讨论：} 模型风险（Model Risk）——每个金融模型都是对现实的高度简化，模型失效不是因为"错了"，而是因为"在关键的时候它不再成立"。LTCM的教训：第一，杠杆放大了一切（50倍杠杆意味着2\%的价格反向变动就清零了资本金）；第二，尾部风险不能被"历史数据"估计——过去100年都没有发生的事，明天可能发生。

\textbf{练习：} 给定一个收敛交易策略（做空高估债券/做多低估债券），计算在不同杠杆率（5x, 10x, 20x, 50x）下，多少bp的价差反向变动会导致爆仓。

\textbf{学海拾遗：} Nassim Taleb的《黑天鹅》——2007年出版，恰好写在金融危机前。Taleb的核心论点：人们用"正态分布"来理解极端事件，但市场的极端事件（金融危机）本质上是由一个完全不同的统计分布（幂律分布）产生的。LTCM、2008年危机、2023年硅谷银行——都是"黑天鹅"。

\item[\textbf{第18课}] \textbf{次贷危机中的"做空者"——信用违约互换的另一种用法} \draft

\textbf{情景：} Michael Burry——一个对冲基金经理，在2005年就发现了次贷市场的巨大泡沫。他做空了MBS——通过买入CDS（信用违约互换）来赌房贷违约。他的基金承受了多年的浮亏和投资人的撤资——然后在2008年赚了约7亿美元。这个故事后来被拍成了电影《大空头》。

\textbf{定性：} "做空"的本质——通过借入并卖出、或者通过衍生品（CDS, Put Option）来从价格下跌中获利。做空的社会价值——做空者揭露泡沫、促进价格发现、提供市场流动性。

\textbf{定量：} CDS定价——保费 = 违约概率 × (1 - 回收率)。做空MBS的损益结构——每月支付的CDS保费（固定成本）vs 违约时的赔付（可变收益）。

\textbf{讨论：} 为什么很少人做空？——做空有"无限损失"的风险（理论上股价可以涨到无穷大），而且做空者面临"轧空"（Short Squeeze）——当被做空的股价快速上涨时，做空者被迫买入平仓，进一步推高股价。2021年GameStop事件就是一次史诗级的轧空——散户在Reddit上联合逼空了Melvin Capital，使其亏损了约70亿美元。

\textbf{练习：} 计算做空一只股票的情景分析：当前股价50元，你以50元借入1000股卖出，得到5万元。如果股价涨到100元/200元/跌到10元，你的盈亏分别是多少？（含融券成本）

\textbf{学海拾遗：} 为什么不禁止做空？——中国证监会2023年曾出台限制做空的措施。学术研究表明：限制做空通常会降低市场效率（价格发现变慢）、增加波动性（因为坏消息只能通过卖股票来释放）、损害流动性。做空不是危机的原因——它是危机预警的信号弹。

\item[\textbf{第19课}] \textbf{2023年硅谷银行——数字时代的银行挤兑} \draft

\textbf{情景：} 2023年3月9日，硅谷银行（SVB）的储户在几个小时内通过网上银行系统取走了420亿美元——约占该银行存款总额的25\%。第二天SVB被FDIC接管。这是历史上第一次完全由社交媒体触发的银行挤兑。

\textbf{定性：} SVB的资产负债表错配——负债端是随时可取的创业公司存款（无保险存款占比94\%），资产端是持有到期的长期国债和MBS。美联储加息500bp后，这些"持有到期"资产的市值蒸发约180亿美元。

\textbf{定量：} 持有到期（HTM）vs 可供出售（AFS）的分类会计——HTM资产不需要按市值计价（Mark-to-Market），因此180亿美元的亏损被藏在资产负债表里没有被公开。但当银行被迫卖资产来应对挤兑时——HTM转AFS——亏损瞬间暴露。

\textbf{讨论：} U.S. Treasury Secretary Janet Yellen："SVB的倒闭不是因为资不抵债——它实际还有偿付能力（按持有到期的账面价值）。SVB倒闭纯粹是因为流动性危机——资产有正的价值，但无法在短时间内以不亏损的价格卖出。"这个分析准确吗？实际上，如果SVB能从美联储借款到所有债券到期——它不会破产。但市场没有给它这个时间。

\textbf{练习：} 计算SVB的资产端久期缺口——资产久期约6.2年，负债久期约1.2年。当利率上升500bp时，净资产的变化是多少？

\textbf{学海拾遗：} 银行业的会计问题——HTM分类让银行可以不按市价计量债券持仓价值。这导致了一个荒谬的情况：SVB的资产负债表上，同一批国债被注明是"持有到期"就值面值，如果改为"可供出售"就亏180亿。数字没有变、风险没有变——只因为分类变了，报表就完全不一样。会计准则是否应该修改以消除这种"会计错觉"？

\item[\textbf{第20课}] \textbf{GameStop、Meme股票与散户化交易的制度冲击} \draft

\textbf{情景：} 2021年1月，WallStreetBets（Reddit上的一个投资社区）的散户们发现GameStop（一只被对冲基金大量做空的股票）——他们集体买入看涨期权和现货，导致股价在两周内从20美元涨到483美元。做空GameStop的Melvin Capital亏损约70亿美元后被关闭。

\textbf{定性：} 轧空（Short Squeeze）的机制——被大量做空的股票价格上涨→做空者被迫买入平仓→进一步推高价格→恶性循环。Gamma Squeeze的机制——大量买入看涨期权→做市商对冲Delta时买入现货→现货价格上涨→更多期权变为实值。

\textbf{定量：} 计算空头比率（Short Interest Ratio）——GameStop空头比率一度超过140\%（因为同一股票被多次借出做空）。当价格从20美元涨到483美元时，做空者的总损失约300亿美元。

\textbf{讨论：} GameStop事件的核心问题——"市场是有效率的"这一信条受到了挑战。一个基本面完全恶化（实体店、疫情期间收入跳水）的公司的股价可以被社交媒体推高24倍。这是"市场失灵"还是"市场在起作用"？从不同的利益相关者角度看，答案不同。

\textbf{练习：} 计算GameStop轧空场景中的做空者累计浮亏。假设做空平均成本20美元、融券费用年化5\%，股价最高到483美元时，1000万股空头的浮亏是多少？

\textbf{学海拾遗：} 交易基础设施的变革——零佣金交易（Robinhood）、社交交易（Reddit, Discord）、期权散户化。这里面最危险的制度创新是零保证金期权交易——让完全没有风险管理知识的零售投资者交易极其高杠杆的产品。2024年美国SEC正在考虑限制零售期权交易的杠杆水平——但这个"潘多拉的盒子"已经打开了。

\item[\textbf{第21课}] \textbf{中国金融风险全景图} \draft

\textbf{情景：} 2020年恒大债务危机、2023年碧桂园违约、中植系暴雷（2023）、村镇银行风险（河南村镇银行2022）——中国金融体系面临的风险是结构性的还是周期性的？

\textbf{定性：} 中国金融风险的四大领域：
\begin{enumerate}
  \item \textbf{房地产}——开发商债务约4万亿人民币（含离岸债），其中大量是"明股实债"的不透明债务结构。恒大总负债2.4万亿——是中国历史上最大的企业债务违约。
  \item \textbf{地方政府隐性债务}——估计约40-60万亿。通过城投公司借的债——本质上是地方政府担保的、但没有在地方政府资产负债表上列出来的债务。
  \item \textbf{影子银行}——2017-2024年从100万亿压缩到约60万亿，但结构型产品仍然大量存在。
  \item \textbf{中小银行风险}——部分城商行和农商行的不良贷款率被低估，拨备覆盖不足。
\end{enumerate}

\textbf{定量：} 中国非金融部门总债务/GDP约280\%——与发达国家相当，但问题是地方政府和企业债务占比过高。如果GDP增速从8\%降至4\%——债务可持续性会迅速恶化。

\textbf{讨论：} "中国版雷曼时刻"会发生吗？一个常见的分析框架：中国金融风险的特点是"高且集中"但不一定是"系统性互连"——因为国有银行和国有企业之间的风险由中央政府信用兜底。但"道德风险"和"监管俘获"是长期隐患。

\textbf{练习：} 分析恒大的简化资产负债表——资产端（土地储备/在建项目/现金）vs 负债端（银行贷款/债券/应付账款/应付税费）。当房价下跌30\%时，净资产的变动。

\textbf{学海拾遗：} 中国的地方债问题实际上是"中央—地方财政关系"的问题——1994年分税制改革后，地方政府承担了约85\%的公共服务支出但只获得约50\%的财政收入。土地财政替代了正式的"地方税"——当房地产市场下行，这个"替代性税基"消失了，地方债风险暴露。出路：建立地方税体系（房产税、消费税分成）和硬化地方政府预算约束。

\item[\textbf{第22课}] \textbf{数字货币与DeFi——金融体系的范式转移？} \draft

\textbf{情景：} 2021年11月，比特币价格达到69000美元——一个由算法创造的、没有任何物理形态的"货币"，总市值超过1.2万亿美元。2022年11月，FTX——全球第二大加密货币交易所——在几天内崩溃，客户损失约80亿美元。加密世界经历了从"革命"到"闹剧"的完整周期。

\textbf{定性：} 区块链的核心思想——去中心化（Decentralization）、不可篡改（Immutability）、无需信任（Trustless）。PoW vs PoS。智能合约——在区块链上运行的自动执行程序。

\textbf{定量：} 比特币的总量上限2100万枚——算法决定的稀缺性。以太坊的EIP-1559机制——销毁手续费base fee，让ETH可能变成通缩货币（实际：2021-2024年ETH供应量从1.2亿降到1.2亿——基本不变）。

\textbf{讨论：} 加密货币是"货币"还是"资产"？——美国SEC认为除了比特币之外的大多数加密资产是"证券"（应该被监管），而加密社区认为它们是"商品"。这个分类决定了由谁监管——以及未来这个行业的生存空间。

\textbf{练习：} 计算一个比特币挖矿的盈亏平衡点——假设矿机成本、电费、全网算力、区块奖励（2024年每个区块3.125 BTC），给定比特币价格，算力成本vs收益。

\textbf{学海拾遗：} 数字人民币（e-CNY）——全球最大的央行数字货币试点项目。与比特币不同：e-CNY是完全中心化的、由央行发行的法律货币的数字版本。它的设计目标不是"颠覆"现有金融体系，而是"替代"流通中的现金（M0）——降低现金运营成本、增强反洗钱能力、在无网络环境下的可离线支付。到2024年，e-CNY累计交易额超过2万亿人民币。但技术和商业上的挑战：用户习惯、商户受理率、与支付宝/微信支付的竞争。

\item[\textbf{第23课}] \textbf{量化交易与高频交易（HFT）——算法操纵市场了吗？} \draft

\textbf{情景：} 2010年5月6日，美国股市发生了"闪电崩盘"（Flash Crash）——道琼斯指数在几分钟内暴跌约1000点（约9\%），然后同样快速地恢复了。原因：一个交易员通过算法卖出了41亿美元的E-Mini S\&P 500期货合约——然后高频交易算法相互影响，造成了"流动性黑洞"。

\textbf{定性：} 量化交易的发展——从统计套利到高频做市到机器学习策略。做市商（Market Maker）的作用——提供流动性、缩小买卖价差。高频交易（HFT）的策略分类：做市、套利、动量点火（Momentum Ignition）、结构性产品套利。

\textbf{定量：} 闪崩的量化分析——订单薄失衡指标（Order Imbalance）：当买单和卖单的比例差距超过某一个阈值时，做市商退出市场，导致瞬间流动性枯竭。HFT的社会价值——学术研究表明HFT降低了买卖价差约50-70\%（降低了交易成本），但也增加了市场脆弱性（闪崩风险）。

\textbf{讨论：} 量化交易是"公平"的吗？——HFT比人类快毫秒级别，这是否构成"不公平竞争"？公平的竞争环境需要所有人有同样的信息获取速度——而HFT的竞争本质上就是速度竞争。2014年Michael Lewis的《Flash Boys》引发了关于HFT的大辩论。

\textbf{练习：} 模拟一个简单的做市策略——给定买卖价差和库存成本，计算在给定波动率下做市商的最优买卖报价宽度。

\textbf{学海拾遗：} 中国的量化交易——2024年中国量化私募管理规模约1.5万亿。2024年2月中国股市大跌，量化产品的DMA（多空收益互换）策略大规模平仓，加剧了市场下跌。中国证监会随后加强了对量化交易的管制——包括申报制度、程序化交易报备、高频交易特别监管。量化在A股的表现因子与美股有系统性差异——小盘因子在中国更强、动量因子更短——这与A股的散户主导结构有关。

\item[\textbf{第24课}] \textbf{金融伦理——资本市场到底在服务谁？} \draft

\textbf{情景：} 2015年A股股灾期间，一些上市公司的大股东通过"高送转"（每10股送20股）的消息拉高股价后减持——市值蒸发了几十万亿，散户投资者损失惨重。这些行为在技术上是"合规"的——但它道德吗？

\textbf{定性：} 金融伦理的核心框架——投资者保护 vs 资本形成 vs 市场效率。利益冲突的四种常见形式：委托-代理问题、内部交易、市场操纵、利益输送。

\textbf{定量：} 内部交易的实证——研究发现内部交易人在交易后的6个月内获得了约3-5\%的超额收益（超过市场平均）。虽然3-5\%看似小，但乘以巨大的交易量——几十亿的不义之财。

\textbf{讨论：} 金融市场的核心矛盾：它应该服务"谁"？——是服务"资本的需求方"（企业融资，降低资本成本），还是服务"资本的供给方"（投资者保护，公平交易），还是服务"整个社会的公共利益"（资本配置效率、经济增长）？这三个目标在理论上可以统一，但在实践中经常冲突——而且不同国家的答案截然不同。华尔街的灵魂拷问：如果金融创新使银行高管获得了1000万美元奖金，但同时增加了金融危机让全社会承担的成本——这是创造了价值，还是转移了风险？

\textbf{练习：} 案例分析——中概股瑞幸咖啡（Luckin Coffee）财务造假事件（2020）。瑞幸通过伪造交易额（约22亿人民币）来虚增收入，被做空机构浑水（Muddy Waters）揭露后股价暴跌85\%以上。

\textbf{学海拾遗：} 绿色金融（Green Finance）与ESG投资——全球ESG基金管理的资产在2021年达到约3万亿美元（2024年略有回落）。ESG评分能否预测超额收益？结论：不能——但ESG评分低（有重大争议）的公司明显更可能有负面事件（罚款、诉讼、声誉损失）。中国ESG的特色——"共同富裕"和"双碳"目标被纳入ESG框架，形成了中国特色。

\end{enumerate}

\newpage

\section*{总统计}

\begin{center}
\begin{tabular}{ll}
\toprule
\textbf{卷} & \textbf{课数} \\
\midrule
第一卷：金融学基础 & 8课（第1-8课）\\
第二卷：衍生品与风险管理 & 6课（第9-14课）\\
第三卷：金融创新与系统性风险 & 10课（第15-24课）\\
\midrule
总计 & 24课 \\
\bottomrule
\end{tabular}
\end{center}

\vspace{0.3cm}

\noindent\textbf{前置要求：}
\begin{itemize}
  \item 第一卷：Algebra II（含指数/对数运算）+ 基础概率
  \item 第二卷：第一卷完成 + 基础微积分（导数和积分的基本概念）
  \item 第三卷：第一卷+第二卷完成
\end{itemize}

\noindent\textbf{每课通用结构：}
\begin{quote}
情景引入（一个真实的事件/场景）→ 定性分析（直觉和框架）→ 定量分析（数学模型和计算）→ 总结（核心结论）→ 讨论题（开放性问题）→ 课后练习（计算题+分析题）→ 学海拾遗（学术细节/历史花絮/扩展阅读）
\end{quote}

\noindent\textbf{当前完成状态：} 第一卷8课中第1-6、8课有初稿（来源于现有 classroom\_materials.tex），第7课和全部第二卷、第三卷为待写大纲。

\vspace{0.5cm}
\noindent\textcolor{gray}{\small 大纲版本：2026-05-23 · 审核通过}
\end{document}
